% =============================================================================
%  SLIDE 4 – Position of Alternate Interior Angles (4 frames)
%    1: title only                        2: parallel lines + transversal
%    3: pair 1 (A1, B1)                   4: pair 2 (A2, B2)
% =============================================================================
\begin{frame}
\begin{tikzpicture}[remember picture,overlay,shift={(current page.south west)},x=1cm,y=1cm]

  \fill[white] (0,0) rectangle (16,9);

  % ---- header ----------------------------------------------------------------
  \node[anchor=west,text=NavyDeep] at (0.8,8.35)
       {{\fontsize{12}{17}\selectfont\muol{រំលឹកមេរៀនចាស់}}\ \ |\ \ %
        {\fontsize{10.5}{15}\selectfont\bfseries ទីតាំង និងលក្ខណៈមុំឆ្លាស់ក្នុង}};
  \draw[NavyDeep!18,line width=0.6pt] (0.8,7.98)--(15.2,7.98);

  % ---- figure ----------------------------------------------------------------
  \only<2->{\altbase}

  % ---- pair 1 : A1 (lower-left of A) and B1 (upper-right of B) ---------------
  \only<3->{
    \angmark{A}{180}{240}{1.05}{PairOne}{0}{$A_1$}{0.62}
    \angmark{B}{0}{60}{1.05}{PairOne}{0}{$B_1$}{0.62}
    \callout{cA1}{4.3,5.0}{$(A)+(210:1.17)$}{PairOne}
    \callout{cB1}{11.2,3.6}{$(B)+(30:1.17)$}{PairOne}
    \node[anchor=west,text=PairOne] at (1.4,1.95)
         {\fontsize{13}{17}\selectfont\bfseries គូទី១ $\angle A_1$ និង $\angle B_1$};
  }

  % ---- pair 2 : A2 (lower-right of A) and B2 (upper-left of B) ---------------
  \only<4->{
    \angmark{A}{240}{360}{1.05}{PairTwo}{0}{$A_2$}{0.62}
    \angmark{B}{60}{180}{1.05}{PairTwo}{0}{$B_2$}{0.62}
    \callout{cA2}{11.6,5.0}{$(A)+(300:1.17)$}{PairTwo}
    \callout{cB2}{4.3,3.9}{$(B)+(120:1.17)$}{PairTwo}
    \node[anchor=west,text=PairTwo!85!black] at (9.4,1.95)
         {\fontsize{13}{17}\selectfont\bfseries គូទី២ $\angle A_2$ និង $\angle B_2$};
  }

  % ---- footer ----------------------------------------------------------------
  \fill[NavyDeep] (0,0) rectangle (16,0.65);
  \shade[shading=axis,left color=GoldAccent,right color=GoldAccent!0] (0,0.65) rectangle (16,0.71);
  \node[anchor=west,inner sep=0,text=white] at (0.8,0.28)
       {\fontsize{8.5}{11}\selectfont \TeacherName\ \ |\ \ \color{SkyLight}គណិតវិទ្យា ថ្នាក់ទី៨};
  \node[anchor=east,inner sep=0,text=GoldAccent] at (15.2,0.28)
       {\fontsize{8.5}{11}\selectfont\muol{មេរៀនទី១២}\ \ \color{white}ការប្រៀបធៀបត្រីកោណ};

\end{tikzpicture}
\end{frame}
